\section*{Bab \ref{ch:g10:concentration-and-dilution} --- Konsentrasi dan Pengenceran}

\begin{solution}{exo:g10:concentration-and-dilution:1}
$C_m = \qty{5.0}{g} / \qty{0.250}{L} = \qty{20}{g/L}$.
\end{solution}

\begin{solution}{exo:g10:concentration-and-dilution:2}
$c = \qty{0.050}{mol} / \qty{0.200}{L} = \qty{0.25}{mol/L}$.
\end{solution}

\begin{solution}{exo:g10:concentration-and-dilution:3}
$M(\ce{CuSO4}) = 63.5 + 32.1 + 4 \times 16.0 = \qty{159.6}{g/mol}$;
$n = 0.10 \times 0.1000 = \qty{0.0100}{mol}$; $m = 0.0100 \times 159.6 =
\qty{1.60}{g}$.
\end{solution}

\begin{solution}{exo:g10:concentration-and-dilution:4}
$F = 1.0 / 0.050 = 20$; $V_0 = \qty{100.0}{mL} / 20 = \qty{5.0}{mL}$.
\end{solution}

\begin{solution}{exo:g10:concentration-and-dilution:5}
$C_m = c \times M = 0.10 \times 58.5 = \qty{5.85}{g/L}$, sekitar
\qty{5.9}{g/L}.
\end{solution}

\begin{solution}{exo:g10:concentration-and-dilution:6}
$V_0 = c_1 V_1 / c_0 = 0.020 \times 250.0 / 0.50 = \qty{10.0}{mL}$
(\omterm{def:g10:concentration-and-dilution:dilution}{faktor pengenceran} 25). Sebuah pipet volume \qty{10.0}{mL} dan sebuah
labu ukur \qty{250.0}{mL}.
\end{solution}

\begin{solution}{exo:g10:concentration-and-dilution:7}
Labu ukur memiliki satu tanda batas saja, dibuat untuk satu volume dan
teliti sampai sebagian kecil persen; garis skala pada gelas kimia hanya
petunjuk kasar. Demikian pula, pipet volume mengalirkan satu volume
dengan sangat teliti, sedangkan gelas ukur lebih lebar dan dibaca kurang
teliti.
\end{solution}

\begin{solution}{exo:g10:concentration-and-dilution:8}
Di antara \qty{0.02}{mol/L} dan \qty{0.04}{mol/L}: sekitar
\qty{0.03}{mol/L}. Agar lebih teliti, tambahkan \omterm{def:g10:concentration-and-dilution:standard}{larutan standar} di antara
keduanya, misalnya pada 0.025, 0.030, dan \qty{0.035}{mol/L}: skala yang
lebih rapat di tempat yang diperlukan.
\end{solution}

\begin{solution}{exo:g10:concentration-and-dilution:9}
$M(\ce{C6H12O6}) = \qty{180.0}{g/mol}$; $c = 50 / 180.0 =
\qty{0.28}{mol/L}$.
\end{solution}

\begin{solution}{exo:g10:concentration-and-dilution:10}
$M = \qty{342.0}{g/mol}$, jadi $n = \qty{1.00}{mol}$ dan $c =
\qty{1.00}{mol/L}$. Setelah diencerkan sepuluh kali: $c = \qty{0.100}{mol/L}$, yaitu
$C_m = 0.100 \times 342.0 = \qty{34.2}{g/L}$. Satu sendok
\qty{0.020}{L} memuat $34.2 \times 0.020 = \qty{0.68}{g}$ sukrosa.
\end{solution}

\begin{solution}{exo:g10:concentration-and-dilution:11}
(a) Massa \omterm{def:g6:solutions-and-solubility:solute}{zat terlarut} yang sama dalam volume dua kali lipat:
\qty{6.0}{g/L}. (b) Volumenya kira-kira menjadi setengah, massa \omterm{def:g6:solutions-and-solubility:solute}{zat
terlarutnya} tidak berubah: sekitar \qty{24}{g/L}.
\end{solution}

\begin{solution}{exo:g10:concentration-and-dilution:12}
\begin{enumerate}
  \item Tiga \omterm{def:g10:concentration-and-dilution:dilution}{pengenceran} 10 kali menghasilkan \omterm{def:g10:concentration-and-dilution:dilution}{pengenceran} $10^3$ kali:
        $\qty{0.10}{mol/L} / 1000 = \qty{1.0e-4}{mol/L}$.
  \item Satu \omterm{def:g10:concentration-and-dilution:dilution}{pengenceran} 1000 kali akan memerlukan volume \omterm{def:g10:concentration-and-dilution:dilution}{larutan induk}
        yang sangat kecil (\qty{0.1}{mL} ke dalam \qty{100}{mL}), yang
        mustahil diukur dengan teliti memakai pipet, atau labu yang sangat
        besar (\qty{1}{mL} ke dalam \qty{1}{L}). Tiga \omterm{def:g10:concentration-and-dilution:dilution}{pengenceran} sepuluh
        kali memakai pipet dan labu yang biasa.
\end{enumerate}
\end{solution}

\begin{solution}{exo:g10:concentration-and-dilution:13}
\begin{enumerate}
  \item $M(\ce{CaCl2}) = 40.1 + 2 \times 35.5 = \qty{111.1}{g/mol}$;
        $n = 11.1 / 111.1 = \qty{0.100}{mol}$; $c = 0.100 / 0.5000 =
        \qty{0.200}{mol/L}$.
  \item Setiap \ce{CaCl2} menghasilkan satu \ce{Ca^2+} dan dua \ce{Cl-}:
        $[\ce{Ca^2+}] = \qty{0.200}{mol/L}$ dan $[\ce{Cl-}] =
        \qty{0.400}{mol/L}$.
\end{enumerate}
\end{solution}

\begin{solution}{exo:g10:concentration-and-dilution:14}
\begin{enumerate}
  \item $r = \qty{0.60}{cm}$, $h = \qty{0.20}{cm}$:
        $V = \pi \times 0.60^2 \times 0.20 = \qty{0.23}{cm^3}$, yaitu
        \qty{0.23}{mL} air yang berlebih.
  \item \omterm{def:g6:solutions-and-solubility:solute}{Zat terlarut} yang sama berada dalam \qty{100.23}{mL} alih-alih
        \qty{100.00}{mL}: konsentrasinya terlalu rendah sebesar
        $0.23 / 100.23 \approx
        \qty{0.23}{\percent}$.
\end{enumerate}
\end{solution}

\begin{solution}{exo:g10:concentration-and-dilution:15}
$n = 0.20 \times 0.100 + 0.10 \times 0.300 = 0.020 + 0.030 =
\qty{0.050}{mol}$ dalam \qty{0.400}{L}: $c = \qty{0.125}{mol/L}$. Hasil
itu bukan rata-rata biasa, \qty{0.15}{mol/L}: larutan yang lebih encer
tiga kali lebih banyak, sehingga \omterm{def:g6:pure-substances-and-mixtures:mixture}{campurannya} lebih dekat ke
\qty{0.10}{mol/L}. Hasil itu adalah rata-rata yang dibobot dengan
volume.
\end{solution}

\begin{solution}{pb:g10:concentration-and-dilution:1}
\textbf{1.} \qty{0.9}{g} per \qty{0.100}{L}: $C_m = \qty{9.0}{g/L}$.

\textbf{2.} \qty{9.0}{g} per \qty{1000}{mL}, yaitu \qty{9.0}{mg} per
mililiter.

\textbf{3.} $m = 9.0 \times 0.500 = \qty{4.5}{g}$.

\textbf{4.} Ya: sekitar \qty{9}{g} per liter dibandingkan sekitar
\qty{360}{g} per kilogram air pada keadaan jenuh, kira-kira empat puluh
kali lebih kecil.

\textbf{5.} $9.0 \times 1.0 = \qty{9.0}{g}$ garam.

\textbf{6.} $M(\ce{NaCl}) = 23.0 + 35.5 = \qty{58.5}{g/mol}$.

\textbf{7.} $c = 9.0 / 58.5 = \qty{0.154}{mol/L}$.

\textbf{8.} $n = 0.154 \times 0.500 = \qty{0.0769}{mol}$ (atau
$4.5 / 58.5$).

\textbf{9.} Setiap \ce{NaCl} menghasilkan satu \ce{Na+} dan satu
\ce{Cl-}: keduanya \qty{0.154}{mol/L}.

\textbf{10.} $0.0769 \times \num{6.02e23} = \num{4.63e22}$ \omterm{def:g9:ions:ion}{ion} natrium.

\textbf{11.} $F = 200 / 9.0 = 22.2$.

\textbf{12.} Sebanyak \qty{4.5}{g} garam dalam kantong itu seluruhnya
berasal dari larutan pekat: $V_0 = \qty{4.5}{g} / \qty{200}{g/L} = \qty{0.0225}{L} =
\qty{22.5}{mL}$.

\textbf{13.} $V_0 = V_1 / F = 500 / 22.2 = \qty{22.5}{mL}$.

\textbf{14.} $c_0 = 200 / 58.5 = \qty{3.42}{mol/L}$;
$c_0 V_0 = 3.42 \times 0.0225 = \qty{0.0769}{mol}$ dan
$c_1 V_1 = 0.154 \times 0.500 = \qty{0.0769}{mol}$: sama.

\textbf{15.} Sekitar $500 - 22.5 = \qty{478}{mL}$ air steril.

\textbf{16.} Pipet ukur \qty{25}{mL} (atau buret), yang dibaca sampai
\qty{0.1}{mL}.

\textbf{17.} Di dalam labu ukur \qty{500.0}{mL}, dicukupkan sampai tanda
batas.

\textbf{18.} $C_m = 200 \times 23.0 / 500 = \qty{9.20}{g/L}$, terlalu
pekat sebesar $0.20 / 9.0 \approx \qty{2.2}{\percent}$.

\textbf{19.} \qty{22.5}{mL} larutan pekat (``20 \%'') untuk satu kantong
\qty{500}{mL}.
\end{solution}
